% =============================================================================
% PLANTILLA DE TESIS CON CITACIÓN IEEE
% =============================================================================
\documentclass[12pt, a4paper, oneside]{report}

% --- Codificación e idioma ---
\usepackage[utf8]{inputenc}
\usepackage[T1]{fontenc}
\usepackage[spanish,es-tabla]{babel}
\usepackage{lmodern} % Fuente vectorial

% --- Márgenes (ajusta según tu universidad) ---
\usepackage[left=3cm, right=2.5cm, top=3cm, bottom=2.5cm]{geometry}

% --- Matemáticas y símbolos ---
\usepackage{amsmath, amssymb, amsthm}
\usepackage{siunitx} % Unidades del SI

% --- Gráficos y figuras ---
\usepackage{graphicx}
\usepackage{float}
\usepackage{subcaption} % Subfiguras

% --- Tablas profesionales ---
\usepackage{booktabs}
\usepackage{multirow}
\usepackage{array}

% --- Espaciado e interlineado ---
\usepackage{setspace}
\onehalfspacing % Interlineado 1.5 (ajusta a \doublespacing si te lo exigen)

% --- Encabezados y pies de página ---
\usepackage{fancyhdr}
\pagestyle{fancy}
\fancyhf{}
\rhead{\thepage}
\lhead{\leftmark}
\renewcommand{\headrulewidth}{0.4pt}

% --- Títulos de capítulos y secciones ---
\usepackage{titlesec}
\titleformat{\chapter}[display]
  {\normalfont\huge\bfseries}{\chaptertitlename\ \thechapter}{20pt}{\Huge}
% Formato del capítulo: "Capítulo 1: Nombre" en una sola línea
\titleformat{\chapter}[block]
  {\normalfont\huge\bfseries}   % Formato general
  {\chaptername\ \thechapter:}  % Etiqueta: "Capítulo 1:"
  {1em}                         % Separación entre etiqueta y título
  {}                            % Código antes del título

% Ajuste de espaciado (opcional, para que no quede muy pegado)
\titlespacing*{\chapter}{0pt}{-20pt}{30pt}


\usepackage[titles]{tocloft}

\renewcommand{\cftchappresnum}{\chaptername~}
\renewcommand{\cftchapaftersnum}{:}
\setlength{\cftchapnumwidth}{6.5em} % Ajusta el ancho según tu numeración


% --- Hipervínculos (DEBE ir al final, casi siempre) ---
\usepackage[colorlinks=true, linkcolor=blue, citecolor=blue, urlcolor=blue]{hyperref}

% --- Numeracion---
\setcounter{secnumdepth}{3}
\setcounter{tocdepth}{3}

% ---Indices ---
\usepackage[titles]{tocloft}
% --- Estilo del índice tipo Word ---
% Puntos de relleno para capítulos, secciones y subsecciones
\renewcommand{\cftchapleader}{\cftdotfill{\cftdotsep}}
\renewcommand{\cftsecleader}{\cftdotfill{\cftdotsep}}
\renewcommand{\cftsubsecleader}{\cftdotfill{\cftdotsep}}

% Separación de los puntos (menor valor = puntos más juntos)
\renewcommand{\cftdotsep}{0.5}

% Ancho reservado para el número de página y margen derecho
% Esto evita que los puntos se corten cuando la página tiene 2 o 3 dígitos
\cftsetpnumwidth{2.5em}
\cftsetrmarg{3.5em}

% (Opcional) Que los capítulos aparezcan en negrita en el índice
\renewcommand{\cftchapfont}{\normalfont\bfseries}
\renewcommand{\cftchappagefont}{\normalfont\bfseries}

\renewcommand{\cftchappresnum}{\chaptername~}
\setlength{\cftchapnumwidth}{6em} % Ajusta el ancho para que quepa "Capítulo 1"
% =============================================================================
% CONFIGURACIÓN DE BIBLIOGRAFÍA CON IEEE (biblatex + biber)
% =============================================================================
\usepackage[
    backend=biber,      % Motor de procesamiento (más moderno que BibTeX)
    style=ieee,         % Estilo de citación IEEE
    sorting=none,       % Ordena las referencias por orden de aparición
    maxbibnames=6,      % Máximo de autores en la bibliografía
    minbibnames=1,      % Mínimo de autores antes de "et al."
    giveninits=true,    % Iniciales en lugar de nombres completos
    isbn=false,         % Ocultar ISBN (opcional)
    url=true,           % Mostrar URLs
    doi=true            % Mostrar DOI
]{biblatex}
\addbibresource{referencias.bib} % Nombre del archivo .bib

% =============================================================================
% DATOS DE LA PORTADA
% =============================================================================
\title{Título de tu Tesis Doctoral}
\author{Nombre Completo del Autor}
\date{\today} % O especifica el mes y año: \date{Septiembre 2026}

% =============================================================================
% INICIO DEL DOCUMENTO
% =============================================================================
\begin{document}

% --- Portada ---
\begin{titlepage}
    \centering
    {\scshape\LARGE Universidad de Ejemplo \par}
    \vspace{1cm}
    {\scshape\Large Facultad de Ingeniería \par}
    \vspace{1.5cm}
    {\huge\bfseries Título de tu Tesis Doctoral \par}
    \vspace{2cm}
    {\Large\itshape Nombre Completo del Autor \par}
    \vfill
    Director(es): \\ Dr. Nombre del Director
    \vfill
    {\large \today \par}
\end{titlepage}

% --- Resumen ---
\begin{abstract}
    Aquí escribes el resumen de tu tesis, de aproximadamente 200-300 palabras.
    Debe incluir el problema, la metodología, los resultados principales y las
    conclusiones.
\end{abstract}

% --- Agradecimientos (opcional) ---
\chapter*{Agradecimientos}
\addcontentsline{toc}{chapter}{Agradecimientos}
Escribe aquí tus agradecimientos.


\tableofcontents
\listoffigures
\listoftables



% --- Capítulos (modular) ---
\chapter*{Introduccion}
La gestión de la distribución, el transporte y la planificación de rutas son temas de gran importancia que han sido abordados a nivel internacional. Son ramas de la logística donde yacen numerosas oportunidades de mejora y por esta razón ha recibido especial atención motivando un intenso trabajo teórico en el desarrollo de modelos y algoritmos. En este contexto surgen los Problemas de Planificación de Rutas de Vehículos (Vehicle Routing Problems, VRP) [1].
En [2] se introduce el "Truck Dispatching Problem", modelando cómo una flota de camiones homogéneos podría satisfacer la demanda de petróleo de varias estaciones de servicio desde un depósito central y con una distancia mínima recorrida. Cinco años después, en [3] se generaliza este a un problema de optimización lineal que se encuentra comúnmente en el ámbito de la logística y el transporte. El VRP es un problema NP-duro [4] por lo que se consideran muy difíciles de resolver. Los problemas de la vida real son considerablemente más grandes y complejos que los escenarios académicos, por lo que las heurísticas y metaheurísticas son necesarias para obtener buenos resultados. Estos algoritmos no garantizan encontrar la solución óptima, pero sí proporcionan tiempos computacionales razonables. La cantidad de métodos de solución introducidos en la literatura académica ha crecido en las últimas décadas [5].
El Problema de Planificación de Rutas de Vehículos con Múltiples Depósitos (Multi-Depot Vehicle Routing Problems, MDVRP) es una variante del VRP clásico que incorpora varios depósitos con una ubicación predefinida. Cada depósito cuenta con una flota de vehículos de capacidad homogénea que tiene como objetivo satisfacer la demanda de un conjunto de clientes dispersos geográficamente. Cada vehículo debe comenzar y terminar su ruta en el mismo depósito, y cada cliente debe ser visitado una única vez. Teniendo en cuenta la distancia entre la ubicación de los clientes y depósitos, el MDVRP busca asignar un cliente a cada depósito y se construye una ruta para cada vehículo dentro de un depósito para atender dichos clientes. El objetivo al crear la ruta es reducir al mínimo el costo total del recorrido [6].
Durante un largo período de tiempo, los métodos heurísticos basados en la búsqueda tabú han sido los más empleados en la literatura [6]. Sin embargo, este problema genera una gran cantidad de combinaciones que se refleja en los altos tiempos computacionales para obtener la solución. El potencial de las técnicas de análisis de clústeres para la solución de los MDVRP ha sido reconocido por autores como [7], [2], [8] y [9]. El análisis de clústeres estudia la división de objetos en grupos basado en una o varias características. Un grupo se define como una zona con muchos objetos parecidos muy cercanos entre sí, que está separada de otros grupos por zonas donde casi no hay objetos. Esta definición de grupo es bastante acorde a la agrupación de clientes por lo que se propone usar el análisis de clústeres en la solución del MDVRP, específicamente en la fase de asignación de clientes a depósitos [7]. En [10] se describe la técnica 'divide y vencerás', utilizada para reducir la complejidad computacional en el empleo heurísticas para resolver el MDVRP.
Sin embargo, aún existen desafíos en cuanto a la eficiencia computacional debido a los altos tiempos y la calidad de las agrupaciones dependiendo del escenario [7]. En la literatura sobresalen algoritmos como Density-Based Spatial Clustering of Applications with Noise (DBSCAN) que permiten identificar clústeres basados en la densidad local [11]. Por otra parte, Clustering Using Representatives (CURE) identifica variaciones en la forma de los clústeres mediante el uso de puntos representativos [12]. 
En este contexto, la Universidad Tecnológica de la Habana “José Antonio Echeverría” (CUJAE) ha desarrollado la Biblioteca de Heurísticas de Asignación para Problemas de Planificación de Rutas de Vehículos con Múltiples Depósitos (BHAVRP). Este componente está diseñado para resolver la fase de asignación en MDVRP utilizando una arquitectura modular y 20 algoritmos repartidos en heurísticas clásicas de asignación y algoritmos de agrupamiento de Minería de Datos adaptados al contexto de los MDVRP [13]. Esta alternativa gratuita y de código abierto supone una ventaja frente a otras herramientas que no tienen una colección de algoritmos tan extensa o son software comercial. 
Esta biblioteca, en la categoría de algoritmos de agrupamiento, incluye algoritmos particionales, algoritmos jerárquicos y un algoritmo basado en densidad. En la literatura se ha sugerido que ciertas métricas de validación internas como Dunn Index o Solhoutte Referencia podrían tener una correlación con la calidad de las rutas en la posterior fase de planificación. En investigaciones anteriores [14], [15], se realizó un análisis experimental donde se concluyó que no había una correlación positiva significativa entre las métricas internas dentro de BHAVRP y la calidad de las rutas obtenidas. Sin embargo, este estudio fue realizado con poca variedad de algoritmos de agrupamiento y sin considerar una amplia variedad de heurísticas de construcción.
Debido a la situación problemática antes descrita, se identifica como problema a resolver la falta de un estudio exhaustivo que determine si existe alguna métrica interna que correlaciones positivamente con la calidad de las rutas obtenidas en la planificación, impide establecer un criterio de selección de los algoritmos de BHAVRP.
Para dar solución al problema planteado se traza como objetivo general: Desarrollar una nueva versión de BHAVRP que incluya más algoritmos de agrupamiento basados en densidad.
Para dar cumplimiento al objetivo general se definieron los siguientes objetivos específicos con sus respectivas tareas:


El cumplimiento del objetivo general tiene como valor práctico un aporte significativo en la resolución de la fase de asignación del problema de planificación de rutas de vehículos con múltiples depósitos. Los resultados tienen valor científico y tecnológico ya que se pone a disposición de la comunidad científica un componente de software que brinda nuevos algoritmos para la solución de MDVRP.
Para alcanzar los objetivos mencionados, el trabajo se encuentra estructurado en tres capítulos, incluyendo secciones importantes como conclusiones, recomendaciones y referencias bibliográficas. El Capítulo 1: Fundamentos teóricos del problema de planificación de rutas y su solución mediante diferentes métodos donde se presenta el estado del arte del MDVRP, las estrategias de solución existentes, así como un análisis comparativo de bibliotecas. El Capítulo 2: Incorporación de algoritmos a la Biblioteca de Heurísticas de Asignación del Problema de Planificación de Rutas con Múltiples Depósitos describe los algoritmos de agrupamiento de BHAVRP, los algoritmos propuestos detallando el funcionamiento de estos según la literatura y las adaptaciones al MDVRP que hay que realizar. Finalmente, el Capítulo 3: Análisis experimental expone varias pruebas unitarias de las funciones más relevantes de los algoritmos implementados, así como un estudio comparativo entre algoritmos de asignación aplicados sobre ciertas instancias y pruebas estadísticas no paramétricas para determinar su rendimiento y aplicabilidad.

\chapter{Fundamentos teóricos del problema de planificación de rutas y su solución mediante diferentes métodos}{1}
\input{capitulos/02_chapter1}
\input{capitulos/02_chapter2}

% --- Bibliografía con IEEE ---
\printbibliography[heading=bibintoc, title={Referencias Bibliograficas}]

% --- Apéndices (opcional) ---
\appendix
\chapter{Apéndice A: Datos Complementarios}
Contenido del apéndice.

\end{document}